% ---------------------------------------------------------------------------
% 論文中文草稿（NTCIR-19 AI CUP VeriPromiseESG Track 投稿用）
% 以 XeLaTeX 編譯：  xelatex paper_draft_zh_codex.tex
% ---------------------------------------------------------------------------
\documentclass[11pt,a4paper]{ctexart}

% Overleaf：Menu -> Compiler 必須設為 XeLaTeX。
% 字體必須明確指定為繁體字集（TC）。若交給 ctex 自動挑選，它在 Overleaf 上會選到
% 簡體字型，繁體字會變成缺字方框。
\usepackage{xeCJK}
\setCJKmainfont{Noto Serif CJK TC}
\setCJKsansfont{Noto Sans CJK TC}
\setCJKmonofont{Noto Sans Mono CJK TC}

\usepackage{geometry}
\geometry{top=1.7cm,bottom=1.8cm,left=1.85cm,right=1.85cm}
\usepackage{booktabs}
\usepackage{array}
\usepackage{amsmath}
\usepackage[table]{xcolor}
\usepackage{enumitem}
\usepackage{hyperref}
\hypersetup{colorlinks=true,linkcolor=black,urlcolor=blue,citecolor=black}
\usepackage{microtype}
\usepackage{fancyhdr}

\pagestyle{fancy}
\fancyhf{}
\fancyhead[L]{NTCIR-19 中文草稿}
\fancyhead[R]{\thepage}
\renewcommand{\headrulewidth}{0.4pt}
\setlength{\parindent}{2em}
\setlength{\parskip}{0.15em}
\setlist{nosep,leftmargin=2em}

\newcommand{\fn}[1]{\texttt{#1}}
\newcommand{\src}[1]{\textit{數值來源：\texttt{\detokenize{research/output/#1}}}}

\title{\textbf{以子句決定邊界、以段落提供脈絡：繁體中文 ESG 承諾與佐證跨度抽取}}
\author{匿名審稿版本}
\date{}

\begin{document}
\ctexset{figurename={圖},tablename={表},abstractname={摘要}}
\maketitle

\begin{abstract}
VeriPromiseESG4K 將企業永續報告中的承諾驗證建成四個分類欄位，另為每個正例標出
承諾與佐證的文字跨度（span）；然而共享任務只評分類，這兩個覆蓋所有正例的欄位尚未
形成評測任務。本文問一個尚未被正式評測的問題：\textbf{繁體中文 ESG 段落中的承諾與
佐證，應以句子還是子句為抽取單元？}資料分析顯示，僅 24\% 的標註邊界落在句界，加入
逗號與頓號後，子句邊界對齊率為 71\%；相應的可達字元 F1 上限由 0.5414 提高至 0.9748。
我們以同一個中文 RoBERTa 編碼器（encoder）、同一組五折切分與訓練設定，比較句子與子句單元。
三個隨機種子的平均結果顯示，子句模型把承諾跨度的字元 F1 由 0.4857 提高至 0.5889，
佐證跨度由 0.6356 提高至 0.6663；移除段落脈絡時，兩者分別下降 0.0590 與 0.0238，
換用另一個中文編碼器則落在種子變異範圍內。這些結果支持一項具體結論：\textbf{此資料的
有效歸納偏誤（inductive bias）是「以子句作決策、以整段作脈絡」的任務形式化。}
我們同時指出字元 F1 會高估「整段照抄」等過度抽取策略，因此必須連同交集聯集比
（intersection over union, IoU）與完全匹配（exact match）報告。

\noindent\textbf{關鍵詞：}ESG、企業承諾、跨度抽取、子句分類、繁體中文、任務形式化

% 摘要數值來源：span_extraction_feasibility.md；tagger_ablation_promise_string.md；
% tagger_ablation_evidence_string.md
\end{abstract}

\section{緒論}

企業永續報告同時包含未來承諾與支持承諾的既有行動。對稽核者而言，只輸出
\fn{Yes}/\fn{No} 或 \fn{Clear}/\fn{Not Clear} 並不充分；系統還應指出判斷依據位於
何處。可定位的文字證據能讓使用者查核錯誤，也能把分類器的輸出轉成可追溯的審閱介面。

VeriPromiseESG4K\cite{day2026} 含 4,000 個段落，其中 2,000 個為測試資料。資料為四個分類欄位評分，
並另附 \fn{promise\_string} 與 \fn{evidence\_string}：前者摘錄承諾語句，後者摘錄佐證
語句。訓練與驗證資料合計來自 49 份 ESG 報告；在具有相應正標籤的資料中，兩個跨度欄位
皆有標註。共享任務的提交格式沒有這兩欄，因此既有評測無法回答「系統是否找對文字」這個
問題。\src{official_test_scoring.md；step12_distribution_comparison.md；span_coverage.md}

本文把這項未評測資源建成一個獨立研究問題：\textbf{應以什麼文字單元抽取承諾與佐證？}
這不是純粹的前處理選擇。句子單元不能表示句內邊界，也不易表示相隔的多段答案；自由字元
起訖模型雖然彈性最高，卻與資料中大量不連續跨度不相容。繁體中文的逗號與頓號提供第三種
選擇：把整段一次編碼，再對每個子句做二元判斷。

本文的貢獻如下：

\begin{enumerate}
  \item 定義 VeriPromiseESG4K 的承諾與佐證跨度評測，並以標註幾何證明句子是錯誤的
        化約單元，而子句只犧牲少量可達上限。
  \item 提出脈絡感知的子句標註器（context-aware clause tagger），以相同編碼器、
        相同切分與相同超參數對照句子 baseline，並以三個隨機種子與消融實驗
        （ablation）定位增益。
  \item 顯示單一字元 F1 不足以評估此任務：整段照抄可取得不低的分數，卻幾乎從不精確；
        因此提出字元 F1、IoU 與完全匹配並列的報告方式。
\end{enumerate}

本文直接為跨度任務另定評量指標與控制比較，使研究問題、方法與證據都落在「定位哪一段
文字」。

\section{相關研究與研究定位}

\paragraph{從 ML-ESG 到承諾驗證。}
ML-ESG 系列依序研究多語 ESG 議題辨識、影響類型與影響期間\cite{chen2023issue,chen2023impact,chen2024duration}。
其中影響期間的工作指出，主觀任務上的模型表現與標註一致度同時較弱，並主張應公開更多
專家討論與判斷細節。ML-Promise\cite{seki2025} 隨後把企業承諾驗證形式化為承諾辨識、
可行動佐證、承諾--佐證清晰度與驗證時間，涵蓋英、法、中、日、韓語。它也提供部分語言的
跨度標註，並以生成模型與檢索增強生成（retrieval-augmented generation, RAG）產生抽取
baseline；然而各語言釋出的文字與跨度欄位並不
完整一致，且其抽取結果採 ROUGE-L，不能直接充當本文資料與指標下的對照。
\src{mlpromise_inventory.md}

\paragraph{VeriPromiseESG4K。}
Day 等人\cite{day2026} 將同一形式化擴充為繁體中文的 4,000 段資料，所有正式標籤與
\fn{promise\_string}、\fn{evidence\_string} 均由人工標註；發布論文以四個分類任務與
LLM／RAG baseline 為主，未報告跨度評測。該論文特別說明以段落而非單句為資料單位，
因為承諾與佐證可能分布於同段不同句；這直接支持本文「子句決策仍須保留段落脈絡」的設計。
發布論文稱語料取自 50 家公司，但釋出資料規範化公司名稱後為 49 個唯一 ticker；本文報告
可重現的釋出檔案統計，並把這項差異視為文件限制。\src{step12_distribution_comparison.md}

\paragraph{特殊場次的研究形態。}
NTCIR-18 AI CUP 特殊場次的 KAP OCR 強化\cite{hsu2025kap}、跨領域檢索集成
\cite{chao2025ensemble}，以及翻譯式與多語檢索比較\cite{chiu2025translation}，都把競賽
資料當成研究材料，另定符合問題本身的指標、baseline、消融或控制比較；其中翻譯式與多語
檢索更同時設置近似上界、配對檢定、錯誤分析與執行成本。本文沿用這個形態：不提出新的
預訓練模型，而是檢驗一個可反駁的任務形式化，並使用跨度專屬指標作答。

\section{任務定義與資料}

\subsection{兩個跨度任務}

對每個段落 $x$，承諾任務預測 \fn{promise\_string} 的字元集合 $P_p$，佐證任務預測
\fn{evidence\_string} 的字元集合 $P_e$。測試集中，承諾任務有 1,625 筆適用資料，佐證
任務有 1,335 筆；訓練加驗證分別有 1,627 與 1,345 筆。答案可能由一段或多段不相鄰文字
組成，故本文把任務視為逐單元的多標籤選取，而不是只能輸出單一起點與終點的抽取式問答。
\src{clause_tagger_promise_string.md；clause_tagger_evidence_string.md；span_extraction_feasibility.md}

\subsection{評量指標}

令預測字元集合為 $P$、標註集合為 $G$。本文使用三個互補指標：

\begin{align}
\mathrm{CharF1} &= \frac{2|P\cap G|}{|P|+|G|},\\
\mathrm{IoU} &= \frac{|P\cap G|}{|P\cup G|},\\
\mathrm{EM} &= \mathbb{I}[P=G].
\end{align}

字元 F1 衡量覆蓋程度；IoU 對多選文字施加更直接的懲罰；完全匹配要求預測與標註一字
不差。三者不能互相替代。以整段文字作為佐證答案時，字元 F1 為 0.6004，但完全匹配只有
0.4\%。其原因是佐證跨度的中位數已占段落 52.2\%，全選策略用極高召回率掩蓋了低精確率。
因此，本文不接受只靠字元 F1 的優劣判定。\src{span_coverage.md；clause_tagger_evidence_string.md}

\section{為何以子句為單元}

我們先做不涉及模型的上限分析。句子以句號、問號、驚嘆號與分號切分；子句再加入逗號與
頓號。對每個標註跨度，選取與其重疊的完整單元，計算在該化約下可能達到的最佳字元 F1。

\begin{table}[ht]
\centering
\small
\begin{tabular}{@{}lrrr@{}}
\toprule
單元 & 中位長度（字元） & 邊界對齊率 & 可達字元 F1 上限 \\
\midrule
句子 & 66 & 24\% & 0.5414 \\
\rowcolor{green!10}
子句 & 14 & 71\% & \textbf{0.9748} \\
\bottomrule
\end{tabular}
\caption{文字單元與標註邊界。邊界對齊率指標註起訖同時落在單元邊界的比例。
\src{span_extraction_feasibility.md}}
\label{tab:ceiling}
\end{table}

表~\ref{tab:ceiling} 的差距來自標註形態，而非模型能力。92\% 的承諾片段完整落在一個
句子內，但片段長度僅占該句中位數的 0.41；選整句通常會多選一半以上文字。另有約五分之一
至四分之一的答案包含多段文字。子句選取同時容許句內邊界與不連續答案，且上限只比自由
字元抽取低 0.0252。\src{span_extraction_feasibility.md}

\section{方法}

\subsection{脈絡感知子句標註器}

本文使用中文 RoBERTa\cite{devlin2019} 的 \fn{hfl/chinese-roberta-wwm-ext-large}
將整個段落編碼一次。斷詞器（tokeniser）提供每個
token 的字元位置；對屬於同一子句的 token 向量取平均，送入共享的二元分類頭，判斷該
子句是否屬於目標跨度。所有判為正例的子句依原順序串接成輸出。因編碼發生在整段層級，
每個子句的表示可以注意承諾與佐證在其他位置的內容。

訓練採五折交叉驗證（five-fold cross-validation）。每一折以四份訓練、在剩餘一份產生
折外預測（out-of-fold prediction）；決策門檻只用折外預測選定，再由五個模型平均測試
機率。承諾與佐證分開訓練。主要設定重跑三個隨機種子；句子 baseline、脈絡消融與編碼器
消融均使用相同折、訓練輪數、批次大小與學習率，且各自重新用折外資料選門檻。
\src{tagger_ablation_promise_string.md；tagger_ablation_evidence_string.md}

\subsection{比較系統}

主要 baseline 是以\textbf{句子}取代子句的監督式模型；這個對照只改變決策單元，能直接
檢驗表~\ref{tab:ceiling} 的標註幾何是否轉化成實際增益。另報告兩種不需訓練的策略：整段
照抄，以及選取含承諾線索詞的句子。生成式 baseline 使用 Qwen3-8B，分為零樣本
（zero-shot）與四樣本（four-shot），要求模型逐字引用原段落。

\section{結果}

\subsection{子句對句子的控制比較}

\begin{table}[ht]
\centering
\small
\setlength{\tabcolsep}{5pt}
\begin{tabular}{@{}lrrrr@{}}
\toprule
 & \multicolumn{2}{c}{承諾跨度} & \multicolumn{2}{c}{佐證跨度} \\
系統 & 字元 F1 & 完全匹配 & 字元 F1 & 完全匹配 \\
\midrule
整段照抄 & 0.4183 & 5.5\% & 0.6004 & 0.4\% \\
Qwen3-8B 零樣本 & 0.2704 & 1.5\% & 0.5537 & 1.4\% \\
Qwen3-8B 四樣本 & 0.3038 & 9.5\% & 0.4727 & 6.8\% \\
\midrule
句子 baseline & 0.4857 & \textbf{14.7\%} & 0.6356 & 8.7\% \\
\rowcolor{green!10}
子句模型（三種子平均） & \textbf{0.5889} & 14.1\% & \textbf{0.6663} & \textbf{12.3\%} \\
\bottomrule
\end{tabular}
\caption{測試結果。監督式兩列使用同一協定；子句列為三個種子平均。生成式 baseline
與監督式模型在相同測試列上評分。\src{tagger_ablation_promise_string.md；
tagger_ablation_evidence_string.md；llm_span_baseline.md}}
\label{tab:main}
\end{table}

子句單元使承諾字元 F1 增加 0.1032，遠大於三個種子的標準差 0.0033；佐證增加
0.0307，其種子標準差為 0.0010。佐證完全匹配亦由 8.7\% 增至 12.3\%。承諾完全匹配則
沒有改善（14.7\% 對 14.1\%）：較細單元能減少部分多選文字，但也增加一個邊界切錯的
機會。本文因此只主張承諾側的重疊品質改善，不主張其完全匹配改善。
\src{tagger_ablation_promise_string.md；tagger_ablation_evidence_string.md}

生成式 baseline 的主要問題不是語意理解，而是忠實抽取。四樣本條件下，承諾任務仍有
518 個回傳片段不在原段落中；佐證任務有 371 個。即使提示明確要求引用，生成模型仍把
抽取做成改寫。\src{llm_span_baseline.md}

\subsection{消融實驗}

\begin{table}[ht]
\centering
\small
\begin{tabular}{@{}lrr@{}}
\toprule
變體 & 承諾字元 F1 & 佐證字元 F1 \\
\midrule
完整子句模型 & \textbf{0.5889} & \textbf{0.6663} \\
改用句子單元 & 0.4857 & 0.6356 \\
移除跨子句脈絡 & 0.5299 & 0.6425 \\
換用 MacBERT 編碼器 & 0.5823 & 0.6676 \\
\bottomrule
\end{tabular}
\caption{消融結果。完整模型為三種子平均；其他變體維持相同訓練協定並重新選門檻。
\src{tagger_ablation_promise_string.md；tagger_ablation_evidence_string.md}}
\label{tab:ablation}
\end{table}

移除段落脈絡後，承諾與佐證分別下降 0.0590 與 0.0238，顯示「以子句作決策」不等於
「只讀子句」。相反地，換用 MacBERT 後的差異為 $-0.0066$ 與 $+0.0013$；前者位於承諾
任務的兩個標準差界線，後者落在佐證的種子變異範圍內。主要增益來自單元與脈絡，而非挑中
特定預訓練編碼器。

單張 RTX 4090 上，五折訓練承諾模型約 296--297 秒、佐證模型約 245 秒；每千個子句推論
約 0.30--0.31 秒。這些時間包含折外與測試推論，顯示子句化沒有引入生成式解碼成本。
\src{tagger_ablation_promise_string.md；tagger_ablation_evidence_string.md}

\section{討論}

\paragraph{研究意義。}
本文的結果不是「又一個競賽模型」。研究對象是標註表示與模型歸納偏誤：句子化約的上限
先由資料幾何預測，再由只改一項設定的監督式實驗驗證；段落脈絡與編碼器消融進一步區分
增益來源。這個設計也產生可移植的假說：對以逗號承載資訊結構、且答案可能不連續的中文
文件抽取，子句分類應優先於整句選取或單一起訖指標。

\paragraph{跨度不等於品質分類。}
找到佐證位置不保證能判定佐證是否清楚。現有的「先定位再判斷」實驗把完整段落、標註
跨度、預測跨度與隨機等長片段作為分類輸入；然而其正對照在 40 個 bootstrap 種子中只有
7 次排除零，平均區間為 $[-0.0392,+0.0003]$。該設計的靈敏度未被證明，因此本文不把
「跨度無助分類」當成結論，也不把跨度模型包裝成分類改進。\src{test3_seed_stability.md}

\paragraph{評測含意。}
字元 F1 適合給部分正確的跨度連續分數，卻會獎勵過度抽取；完全匹配直觀，卻把少一個
標點與漏掉整個子句都計為零。IoU 位於兩者之間。未來共享任務若納入跨度，應三者並列，
並另外報告空輸出、原文不存在片段與多段答案的比例。

\section{限制}

第一，子句規則依中文標點定義；表格、條列與 OCR 錯誤可能使標點不可靠。第二，完全匹配
仍偏低：承諾 14.1\%、佐證 12.3\%，故目前模型適合提供候選文字，不宜宣稱可取代人工
摘錄。第三，句子 baseline 是必要但不充分的對照；本文尚未與同編碼器的 token-level BIO
標註器比較，後者能表示任意邊界，可能縮小或逆轉子句優勢。第四，完整模型雖重跑三個種子，
句子、脈絡與編碼器變體各只有一次訓練；目前的種子標準差是噪音尺度，不是配對顯著性
檢定。第五，模型只在單一繁體中文資料集上驗證，沒有跨公司年份或跨語言外部測試。第六，
原始跨度以分隔符號串接多段文字，清理後仍有少數片段無法在原段落逐字定位；現行評測程式
會忽略未對齊片段，故正式投稿前應修正對齊並重算全部跨度指標。第七，標註跨度本身沒有
獨立人工複核；邊界差異可能反映多個都合理的摘錄方式，而非模型錯誤。
\src{span_clean_redo.md；tagger_ablation_promise_string.md；tagger_ablation_evidence_string.md}

\section{結論}

VeriPromiseESG4K 的跨度標註揭示一個未被分類評測涵蓋的研究問題：承諾與佐證的自然單元
不是句子，而是保留整段脈絡的子句。標註幾何把可達字元 F1 上限由 0.5414 提高到 0.9748；
控制實驗再把承諾與佐證的實際字元 F1 由 0.4857／0.6356 提高至 0.5889／0.6663。
脈絡消融有明顯代價，編碼器替換則沒有超出種子噪音。因而，本文的核心結論不是模型規模，
而是任務形式化：\textbf{以子句決定邊界，以段落提供語境，並用多個跨度指標防止過度抽取
被單一分數掩蓋。}

% 結論數值來源：span_extraction_feasibility.md；tagger_ablation_promise_string.md；
% tagger_ablation_evidence_string.md

\begin{thebibliography}{9}
\small

\bibitem{day2026}
Min-Yuh Day, Wei-Chen Huang, Hsin-Ting Lu, Wen-Ze Chen, Yu-Han Huang,
Ching-Tai Chen, and Yohei Seki. VeriPromiseESG4K: A Traditional Chinese ESG
Promise Verification Dataset. 2026.

\bibitem{seki2025}
Yohei Seki, Hakusen Shu, Anais Lhuissier, Hanwool Lee, Juyeon Kang,
Min-Yuh Day, and Chung-Chi Chen. ML-Promise: A Multilingual Dataset for
Corporate Promise Verification. EMNLP, pages 20353--20366, 2025.

\bibitem{chen2023issue}
Chung-Chi Chen, Yu-Min Tseng, Juyeon Kang, Anais Lhuissier, Min-Yuh Day,
Teng-Tsai Tu, and Hsin-Hsi Chen. Multi-Lingual ESG Issue Identification.
FinNLP, pages 111--115, 2023.

\bibitem{chen2023impact}
Chung-Chi Chen, Yu-Min Tseng, Juyeon Kang, Anais Lhuissier, Yohei Seki,
Min-Yuh Day, Teng-Tsai Tu, and Hsin-Hsi Chen. Multi-Lingual ESG Impact Type
Identification. FinNLP, 2023.

\bibitem{chen2024duration}
Chung-Chi Chen, Yu-Min Tseng, Juyeon Kang, Anais Lhuissier, Yohei Seki,
Hanwool Lee, Min-Yuh Day, Teng-Tsai Tu, and Hsin-Hsi Chen. Multi-Lingual ESG
Impact Duration Inference. FinNLP, 2024.

\bibitem{hsu2025kap}
Hsin-Ling Hsu and Ping-Sheng Lin. KAP: MLLM-assisted OCR Text Enhancement for
Hybrid Retrieval in Chinese Non-Narrative Documents. Proceedings of FinTech in
AI CUP Special Session, NTCIR-18, 2025.

\bibitem{chao2025ensemble}
Chi-Hsiang Chao, Hsin-Fu Chang, and Po-Yuan Teng. Ensemble-Based
Multi-Specialty Retrieval: Integrating Diverse Similarity Metrics for Enhanced
Question Answering. Proceedings of FinTech in AI CUP Special Session,
NTCIR-18, 2025.

\bibitem{chiu2025translation}
Yi-Ting Chiu and Zong-Han Bai. Translation or Multilingual Retrieval?
Evaluating Cross-Lingual Search Strategies for Traditional Chinese Financial
Documents. Proceedings of FinTech in AI CUP Special Session, NTCIR-18, 2025.

\bibitem{devlin2019}
Jacob Devlin, Ming-Wei Chang, Kenton Lee, and Kristina Toutanova. BERT:
Pre-training of Deep Bidirectional Transformers for Language Understanding.
NAACL-HLT, 2019.

\end{thebibliography}

\end{document}
